\documentclass{letter}
\usepackage{hyperref}
\usepackage{amsmath}
\signature{Dawid Kucharski}
\address{Poznan University of Technology\\Institute of Mechanical Technology\\Poznan, Poland\\\href{mailto:dawid.kucharski@put.poznan.pl}{dawid.kucharski@put.poznan.pl}}
\begin{document}

\begin{letter}{The Editor\\Physica A: Statistical Mechanics and its Applications}

\opening{Dear Editor,}

I am pleased to submit the manuscript entitled ``Calibrating the Brody exponent as a quantitative measure of short-range exclusion in 2D spatial point processes'' for consideration for publication in Physica A.

\textbf{Why Physica A?} This paper develops a calibrated measurement framework for short-range exclusion in 2D spatial point processes---a statistical-physics methodology demonstrated on manufactured surfaces, prime number embeddings, and interferometric profilometry. The work bridges random matrix theory (Brody distribution) with established spatial point-process statistics, providing the correct 2D reference frame and an empirical calibration curve (Spearman $\rho=0.988$) that translates the dimensionless Brody parameter into a physically interpretable exclusion radius. The interdisciplinary scope---spanning surface engineering, number theory, and optical metrology---is well-suited to Physica A's readership.

\textbf{Key contributions:}
\begin{enumerate}
  \item The 2D complete-spatial-randomness baseline is recalibrated to $\beta=0.96\pm0.15$ (30 realisations, density-stable), correcting the inappropriate 1D Poisson reference $\beta=0$. A distinct binary-field CSR baseline ($\beta\approx1.46$ at low fill fraction) is identified, with a decision table provided for practitioners.
  \item An empirical $\beta$--$r_{\text{excl}}/\langle\text{NN}\rangle$ calibration curve (Spearman $\rho=0.988$ with effective hard-core radius) translates the dimensionless Brody parameter into a physically interpretable exclusion measure spanning CSR to near-crystalline regimes.
  \item A sparse-integer control proves that the elevated $\beta=2.15$ for 2D prime embeddings is a genuine arithmetic signal ($\Delta\beta=+0.68$ over random-integer control), not a sparsity artefact.
  \item The Cantor embedding null result ($\beta=1.40<\beta_{\text{Bern}}=1.47$, TOST $p<0.01$) demonstrates that 2D exclusion in arithmetic sequences is embedding-created rather than intrinsic---the framework detects both the presence and the absence of exclusion.
  \item Density-thinning, embedding-robustness, spatial-permutation, and peak-detection sensitivity controls isolate genuine exclusion from confounding effects across 58 manufactured surfaces (10 materials, 10 processes).
\end{enumerate}

\textbf{Rigour and reproducibility:} All analyses are accompanied by bootstrap confidence intervals, AIC-based model selection, Benjamini--Hochberg multiple-comparison correction, power analysis, and formal equivalence testing (TOST). The complete analysis code, derived results, and generated figures are archived in a public GitHub repository (github.com/dawidkucharski/brody-2d-exclusion). Raw surface and interferometric data are available upon request. Every numerical claim in the manuscript is verifiable against the deposited data and scripts.

\textbf{Declarations:} This manuscript is original, has not been published elsewhere, and is not under consideration by another journal. The author declares no competing interests. This is a single-author submission. The manuscript complies with Physica A's formatting and ethical guidelines.

\closing{Respectfully,}

\end{letter}
\end{document}
